\section{Appendix 1: Polynomial Chaos Expansion}
A Polynomial Chaos Expansion (PCE) can model a function $f(\vec{x},t;\vec{\xi})$, with independent variables $\vec{x}$, and $t$, and uncertain inputs $\vec{\xi}$. If $\vec{\xi} \in \mathbb{R}^n$, then $n$ can be defined as the \textit{stochastic dimension}, and $f(\vec{x},t;\vec{\xi})$ can be modelled as:
$$
f(\vec{x}, t; \vec{\xi}) \approx \sum_{k=0}^{\text{N\_basis}} c_k(\vec{x},t) \Psi_k(\vec{\xi})
$$
Where:
$$
\text{N\_basis} = \frac{(n + p)!}{n! + p!}
$$
And $p$ is the \textit{Polynomial Order} hyperparameter.

\parskip = \baselineskip
\noindent $\mathbf{\Psi_k(\vec{\xi})}$ - \textbf{Polynomial Basis}: This term represents a collection of polynomials, and is expressed as:
$$
\Psi_k(\vec{\xi}) = \prod_{j=1}^{n} \Psi_{\alpha_j^{(k)}} (\xi_j)
$$
Recall that $n$ is the stochastic dimension. $\Psi_{\alpha_j^{(k)}} (\xi_j)$ is a member of a specific \textit{family} of polynomials depending on the distribution of $\vec{\xi}$:
\begin{itemize}
  \item If $\vec{\xi}$ is \textit{uniformly} distributed, we draw from the \href{https://en.wikipedia.org/wiki/Legendre_polynomials}{Legendre Polynomal Family}.
  \item If $\vec{\xi}$ is distributed according to a \textit{gaussian} distribution, we use the \href{https://en.wikipedia.org/wiki/Hermite_polynomials}{Hermite Polynomal Family}. 
\end{itemize}

\noindent A polynomial of order $k$ in either of these families can be described by a vector of coefficients ranging from the 0th order coefficients up to the $k$th coefficients. However, when we say a polynomial is of order $k$, we mean that the $k^{\text{th}}$ coefficient is 1, and all others are zero.

\noindent Moreover, $\Psi_{\alpha_j^{(k)}}$ is a member of a particular set of a \textit{multi-index} of polynomials.
A multi-index is a set of sets of polynomials.
In particular, it is the set of all sets, each made up of $j$ polynomials, whose total polynomial order equals $k$.
For example:
$$
\Psi_{\alpha_2^{(k)}} = \left\{ (\Psi^0, \Psi^3), (\Psi^3, \Psi^0), (\Psi^1, \Psi^2), (\Psi^2, \Psi^1), (\Psi^2, \Psi^2) \right\}
$$

\noindent $\mathbf{c_k}$ - \textbf{Deterministic Coefficients}: The goal of training a PCE is to find these coefficients. During training, we know $f$ for a range of inputs. Moreover, from the stochastic dimension (enforced by the dimension of $\vec{\xi}$), and user-specified polynomial order, we can calculate $\Psi_k$. All that is left to do is calculate $c_k$.

\noindent We can solve for $c_k$ using either spectral projection or least squares.
While spectral projection takes advantage of the orthogonal nature of our polynomial basis, it's implentation is out of scope here. We instead use a least squares solve to find coefficients.

\subsection{Vector Form of the PCE}
In vector form, a PCE looks like:
$$
\vec{y} = \mathbf{\Psi}(\vec{\xi}) \vec{c}
$$
\subsubsection{Training}
To train a PCE, we create $\text{N\_samples}$ evaluations of $f$, each with a different realization of $\vec{\xi}^{(i)}$. We then populate $\mathbf{\Psi}$ based on: The distribution of $\vec{\xi}$ (Normal or Uniform), a user specified polynomial order $p$, and the stochastic dimension $n$.
Once $\bm{\Psi}$ is calculated, we solve for $\vec{c}$ with a least squares method.

$$
\begin{bmatrix}
f(\vec{x},t ; \vec{\xi}^{(0)}) \\
f(\vec{x},t ; \vec{\xi}^{(1)}) \\
... \\
... \\
f(\vec{x},t ; \vec{\xi}^{(\text{N\_samples} - 1)}) \\
\end{bmatrix}
=
\begin{bmatrix}
\Psi_0(\vec{\xi}^{(0)}) & \Psi_1(\vec{\xi}^{(0)}) & \cdots & \Psi_{\text{N\_basis}}(\vec{\xi}^{(0)}) \\
\Psi_0(\vec{\xi}^{(1)}) & \cdots & \cdots\\
\cdots \\
\Psi_0(\vec{\xi}^{(\text{N\_samples} - 1)}) & \Psi_1(\vec{\xi}^{(\text{N\_samples} - 1)}) & \cdots & \Psi_{\text{N\_basis} - 1}(\vec{\xi}^{(\text{N\_samples} - 1)}) \\
\end{bmatrix}
\begin{bmatrix}
c_0 \\
\cdots \\
c_{\text{N\_basis}}
\end{bmatrix}
$$

\subsubsection{Prediction}
When presented with a novel $\vec{\xi}$ matrix, we populate $\mathbf{\Psi}$ matrix to match the dimension of $\vec{\xi}$, and multiply it against our trained $\vec{c}$ to get an estimate $\tilde{f}(\vec{x},t ; \vec{\xi}_\text{novel})$.